% sections/subsection1/main.tex
\setcounter{section}{0}
\renewcommand{\thesubsection}{\arabic{subsection}.}

\subsection{ Đặt vấn đề}
\hspace*{2em} Trong thời đại số hóa hiện nay, giao tiếp trực tuyến đã trở thành một phần không thể thiếu trong cuộc sống hàng ngày, từ mạng xã hội, công cụ làm việc như Slack hay Microsoft Teams đến tin nhắn truyền thống trên điện thoại thông minh. Tuy nhiên, sự gia tăng của các mối đe dọa an ninh mạng, bao gồm tin tặc, tổ chức không được ủy quyền và các cuộc tấn công kiểu "người trung gian", đã đặt ra nhu cầu cấp thiết về việc bảo mật dữ liệu người dùng. Các nghiên cứu gần đây, như bài báo "Secure Chat Application with an End-to-End Encryption" (CIST 140), chỉ ra rằng mã hóa đầu cuối là giải pháp quan trọng để bảo vệ quyền riêng tư và toàn vẹn thông tin. Đề tài "Ứng dụng RSA và AES trong phát triển hệ thống nhắn tin bảo mật" được lựa chọn nhằm đáp ứng nhu cầu này, tận dụng sức mạnh của các thuật toán mã hóa tiên tiến để xây dựng một ứng dụng nhắn tin an toàn, hiệu quả và thân thiện với người dùng, đồng thời góp phần nâng cao nhận thức về bảo mật trong giao tiếp số.
\subsection{ Mục tiêu}
\hspace*{2em} Mục tiêu của đề tài này là:
\begin{itemize}
  \item \textbf{Phát triển hệ thống nhắn tin với mã hóa đầu cuối:} 
  Dựa trên mô hình của CIST 140, mục tiêu này tham khảo cách bài báo sử dụng RSA cho trao đổi khóa và AES cho mã hóa tin nhắn để đảm bảo chỉ người gửi và người nhận truy cập được nội dung \cite{carranza2024}.

  \item \textbf{Tích hợp RSA và AES:} 
  Tham chiếu từ phương pháp luận của CIST 140, kết hợp RSA cho khóa bất đối xứng và AES cho khóa đối xứng, như được mô tả trong nghiên cứu của Sabah et al. (2005) \cite{sabah2005}.

  \item \textbf{Xây dựng GUI thân thiện:} 
  Sử dụng tkinter như trong CIST 140 để tạo giao diện dễ sử dụng, tham khảo nguyên tắc cân bằng bảo mật và tiện ích từ Tariq et al. (2023) \cite{tariq2023}.

  \item \textbf{Kiểm tra bảo mật:} Tham khảo phương pháp kiểm tra bằng Wireshark trong CIST 140 để đánh giá hiệu quả mã hóa.

  \item \textbf{Mở rộng tính năng:} Lấy cảm hứng từ xu hướng đa phương tiện trong nghiên cứu của Chatterjee et al. (2018) {\cite{chatterjee2018}} và kết luận của CIST 140 (Section 6) để đề xuất phát triển thêm cho tương lai.
\end{itemize}

\subsection{ Phạm vi nghiên cứu}
\hspace*{2em} Phạm vi nghiên cứu tập trung vào việc thiết kế và phát triển một ứng dụng nhắn tin bảo mật sử dụng Python, tích hợp mã hóa văn bản thông qua sự kết hợp của RSA và AES, kết nối qua socket, và xây dựng giao diện người dùng bằng thư viện tkinter, như được triển khai trong dự án. Tuy nhiên, đề tài có giới hạn ở việc chỉ hỗ trợ nhắn tin văn bản, không bao gồm các tính năng đa phương tiện như hình ảnh hoặc video, và chỉ triển khai trong môi trường cục bộ (localhost).

